%Abstract
\setlength{\headheight}{14pt}
\begin{center}
	{\huge \bf Abstract}
	\line(1,0){430}
\end{center}

%The abstract should contain a summary of the work presented in this report in a single paragraph. The paragraph should provide a general background of the problem, methodology and results.

Large Language Models (LLMs) can answer complex scientific questions, but they still produce hallucinations and incorrect reasoning. Multi-agent debate improves reasoning by allowing multiple AI agents to discuss a question before generating a final answer. However, existing debate systems often rely on majority voting or confidence scores. A confident but incorrect majority can influence a correct minority agent, leading the debate to an incorrect conclusion. Existing methods do not verify whether an agent's claims are supported by external scientific evidence before determining its influence. To mitigate this problem, we propose a Trust-Calibrated Multi-Agent Scientific Deliberation framework that improves the reliability of scientific question answering by dynamically adjusting each agent's influence during the debate. The framework breaks each agent's response into atomic claims, retrieves relevant scientific literature, verifies each claim against the retrieved evidence, and updates a trust score for every agent throughout the debate. The final answer is produced using trust-weighted aggregation instead of majority voting. By grounding trust in external scientific evidence, the proposed framework reduces the influence of unsupported claims and helps preserve correct minority opinions during the debate. We expect the framework to reduce hallucinations, improve evidence-based reasoning, prevent incorrect consensus, and generate more reliable answers for scientific question answering without requiring model fine-tuning.

\clearpage
\setlength{\headheight}{12pt}